\documentclass[11pt,a4paper]{article}
\usepackage[utf8]{inputenc}
\usepackage[T1]{fontenc}
\usepackage[catalan]{babel}
\usepackage{lmodern}
\usepackage[sfdefault,lining]{FiraSans}
\usepackage[a4paper,top=2.3cm,bottom=2.2cm,left=1.8cm,right=1.8cm,headheight=1.3cm,headsep=0.4cm,footskip=1cm]{geometry}
\usepackage{amsmath,amssymb}
\usepackage{icomma}
\usepackage{graphicx}
\usepackage[table]{xcolor}
\usepackage{tikz}
\usepackage{fontawesome5}
\usepackage{ccicons}
\usepackage{fancyhdr}
\usepackage{lastpage}
\usepackage{tasks}
\usepackage{needspace}
\usepackage[most]{tcolorbox}
\usepackage{array}
\usepackage[hidelinks,pdftitle={Solucions de la fitxa d'exercicis de la unitat 3: Proporcionalitat i percentatges. 3r d'ESO},pdfsubject={Matemàtiques 3r ESO},pdfkeywords={percentatges, augments, disminucions, interés compost, repartiments, mescles, moviments, proporcionalitat composta}]{hyperref}

\definecolor{navy}{HTML}{1B2748}
\definecolor{accent}{HTML}{E8603A}
\definecolor{teal}{HTML}{1F9E90}
\definecolor{soft}{HTML}{EEF1F8}
\definecolor{muted}{HTML}{5D6685}

\setlength{\parindent}{0pt}
\setlength{\parskip}{3pt}
\settasks{label=\alph*), label-format=\bfseries\color{navy}, label-width=1.3em, item-indent=1.8em, column-sep=0.6em, after-item-skip=2pt, before-skip=3pt, after-skip=4pt}

% ---------- capçalera i peu ----------
\newcommand{\llicencia}{\href{https://creativecommons.org/licenses/by-nc-sa/4.0/deed.ca}{\ccbyncsa}\ \ Llicència Creative Commons Reconeixement-NoComercial-CompartirIgual 4.0 Internacional (CC BY-NC-SA 4.0)}
\pagestyle{fancy}
\fancyhf{}
\fancyhead[L]{\raisebox{-0.35\height}{\includegraphics[height=0.95cm]{logo-gva.png}}}
\fancyhead[C]{\small\color{muted}Solucions de la fitxa · Unitat 3: Proporcionalitat i percentatges · 3r d'ESO}
\fancyhead[R]{\raisebox{-0.35\height}{\includegraphics[height=0.95cm]{logo-vermellar.png}}}
\fancyfoot[L]{\scriptsize\color{muted}\llicencia}
\fancyfoot[R]{\small\color{muted}\thepage}
\renewcommand{\headrulewidth}{0.4pt}
\renewcommand{\headrule}{\color{muted!40}\hrule height \headrulewidth}
\fancypagestyle{primera}{\fancyhf{}\renewcommand{\headrulewidth}{0pt}
  \fancyfoot[L]{\scriptsize\color{muted}\llicencia}
  \fancyfoot[R]{\small\color{muted}\thepage}}

% ---------- nivells i exercicis ----------
\newcommand{\estrelles}[1]{{\color{accent}\ifcase#1\or\faStar\faStar[regular]\faStar[regular]\or\faStar\faStar\faStar[regular]\or\faStar\faStar\faStar\fi}}
\newcounter{exer}
\newcommand{\ex}[1]{\par\needspace{4\baselineskip}\medskip\stepcounter{exer}\noindent{\color{navy}\bfseries\theexer.}\ \ {\small\estrelles{#1}}\ \ }

\newcommand{\seccio}[2]{\par\needspace{10\baselineskip}\bigskip
  \begin{tcolorbox}[colback=navy,colframe=navy,coltext=white,arc=2mm,boxrule=0pt,left=3mm,right=3mm,top=1.2mm,bottom=1.2mm]
  {\large\bfseries #1\quad #2}
  \end{tcolorbox}\vspace{-2pt}}

% ---------- recta real ----------
\newcommand{\rectabuida}{\begin{tikzpicture}[x=0.42cm,baseline=0]
  \draw[-latex,muted] (-6.7,0)--(6.8,0);
  \foreach \i in {-6,...,6}{\draw[muted] (\i,0.12)--(\i,-0.12) node[below=-1pt,font=\scriptsize,text=muted]{$\i$};}
  \end{tikzpicture}}
% \rectaint{x1}{x2}{estil esquerra}{estil dreta}: estil = ple, buit, inf
\newcommand{\extrem}[2]{\ifx#2p\fill[accent] (#1,0) circle (0.13);\else\ifx#2b\filldraw[fill=white,draw=accent,line width=1.2pt] (#1,0) circle (0.13);\fi\fi}
\newcommand{\rectaint}[4]{\begin{tikzpicture}[x=0.42cm,baseline=0]
  \draw[-latex,muted] (-6.7,0)--(6.8,0);
  \foreach \i in {-6,...,6}{\draw[muted] (\i,0.12)--(\i,-0.12) node[below=-1pt,font=\scriptsize,text=muted]{$\i$};}
  \draw[accent,line width=3pt] (#1,0)--(#2,0);
  \ifx#3i\fill[accent] (#1-0.25,0)--(#1+0.15,0.2)--(#1+0.15,-0.2)--cycle;\else\extrem{#1}{#3}\fi
  \ifx#4i\fill[accent] (#2+0.25,0)--(#2-0.15,0.2)--(#2-0.15,-0.2)--cycle;\else\extrem{#2}{#4}\fi
  \end{tikzpicture}}

\newcolumntype{C}[1]{>{\centering\arraybackslash}p{#1}}

% \res: requadre per al resultat
\newcommand{\res}[1]{\boxed{#1}}

\begin{document}
\thispagestyle{primera}

% ---------- encapçalament ----------
\begin{minipage}[c]{0.25\textwidth}\includegraphics[width=\linewidth]{logo-gva.png}\end{minipage}%
\begin{minipage}[c]{0.5\textwidth}\centering
  {\color{navy}\Large\bfseries Solucions de la fitxa de la unitat 3:\\[2pt] Proporcionalitat i percentatges}\\[5pt]
  {\color{accent}\large\bfseries 3r d'ESO}
\end{minipage}%
\begin{minipage}[c]{0.25\textwidth}\hfill\includegraphics[width=0.85\linewidth]{logo-vermellar.png}\end{minipage}

\vspace{6pt}
{\color{navy}\hrule height 1.2pt}
\vspace{8pt}

\begin{tcolorbox}[colback=soft,colframe=soft,arc=2mm,left=3mm,right=3mm,top=1.5mm,bottom=1.5mm]
\small Intenta resoldre cada exercici abans de mirar-ne la solució. Si no t'ix el mateix resultat, repassa els passos o consulta els apunts interactius de la unitat. Els resultats amb decimals estan arredonits a les centèsimes.
\end{tcolorbox}


% =====================================================
\seccio{1}{Percentatges, fraccions i decimals}

\ex{1}\par\smallskip
{\renewcommand{\arraystretch}{1.75}\centering
\begin{tabular}{|C{4.2cm}|C{4.2cm}|C{4.2cm}|}
\hline
\rowcolor{soft}\textbf{Percentatge} & \textbf{Fracció irreductible} & \textbf{Decimal} \\ \hline
$40\,\%$ & $\frac25$ & $0,4$ \\ \hline
$75\,\%$ & $\frac34$ & $0,75$ \\ \hline
$15\,\%$ & $\frac{3}{20}$ & $0,15$ \\ \hline
$8\,\%$ & $\frac{2}{25}$ & $0,08$ \\ \hline
$35\,\%$ & $\frac{7}{20}$ & $0,35$ \\ \hline
$125\,\%$ & $\frac54$ & $1,25$ \\ \hline
$2,5\,\%$ & $\frac{1}{40}$ & $0,025$ \\ \hline
\end{tabular}\par}

\ex{1}
\begin{tasks}(6)
  \task $32\,\%$
  \task $7\,\%$
  \task $150\,\%$
  \task $60\,\%$
  \task $36\,\%$
  \task $12,5\,\%$
\end{tasks}

\ex{2}
\begin{tasks}(6)
  \task $\frac{3}{25}$
  \task $\frac{13}{20}$
  \task $\frac{1}{125}$
  \task $\frac18$
  \task $\frac75$
  \task $\frac{1}{400}$
\end{tasks}

\ex{2} En percentatge: $\frac25=40\,\%$, \ $0,38=38\,\%$, \ $41\,\%$, \ $\frac{7}{16}=43,75\,\%$, \ $0,\overline{4}=44,\overline{4}\,\%$.
\[\res{0,38<\frac25<41\,\%<\frac{7}{16}<0,\overline{4}}\]

\ex{3} A: $\frac{18}{24}=0,75=\res{75\,\%}$. \ B: $\frac{21}{30}=0,7=\res{70\,\%}$. \ C: $\frac{16}{20}=0,8=\res{80\,\%}$. \quad Ordre: \textbf{C, A, B}.

% =====================================================
\seccio{2}{Càlcul de percentatges}

\ex{1}
\begin{tasks}(6)
  \task $30$
  \task $48$
  \task $18$
  \task $30$
  \task $60$
  \task $10$
\end{tasks}

\ex{1}
\begin{tasks}(6)
  \task $47$
  \task $43$
  \task $30$
  \task $35$
  \task $9$
  \task $12$
\end{tasks}

\ex{2} Xiques: $850\cdot0,46=\res{391}$. \ Xics: $850-391=\res{459}$ \ (el $54\,\%$).

\ex{2} Descompte: $45\cdot0,3=\res{13,50\text{ €}}$. \ Paguem: $45-13,50=\res{31,50\text{ €}}$.

\ex{2} Futbol: $1\,200\cdot0,35=\res{420}$. \ Bàsquet: $1\,200\cdot0,28=\res{336}$. \ Tennis: $1\,200\cdot0,15=\res{180}$.\\
Altres ($22\,\%$): $1\,200-420-336-180=\res{264}$.

\ex{3} Ordinadors: $2\,500\cdot0,4=\res{1\,000\text{ €}}$; \ queden $1\,500$ €. \ Programes: $1\,500\cdot0,3=\res{450\text{ €}}$.\\
Formació: $1\,500-450=\res{1\,050\text{ €}}$, \ que és $\frac{1\,050}{2\,500}=0,42=\res{42\,\%}$ del pressupost.

\ex{3} $30\,\%$ de $500=150$ i $20\,\%$ de $150=\res{30}$. \ $20\,\%$ de $500=100$ i $30\,\%$ de $100=\res{30}$: és el mateix.\\
Equival a $0,2\cdot0,3=0,06$, és a dir, al $\res{6\,\%}$ de $500$.

% =====================================================
\seccio{3}{Tant per cent que representa un nombre respecte d'un altre}

\ex{1}
\begin{tasks}(6)
  \task $25\,\%$
  \task $30\,\%$
  \task $75\,\%$
  \task $150\,\%$
  \task $1,5\,\%$
  \task $20\,\%$
\end{tasks}

\ex{1} Aproven $\frac{21}{25}=0,84=\res{84\,\%}$. \ Suspenen $100\,\%-84\,\%=\res{16\,\%}$.

\ex{2}
\begin{tasks}(4)
  \task $27:0,15=\res{180}$
  \task $18:0,4=\res{45}$
  \task $54:1,2=\res{45}$
  \task $5:0,025=\res{200}$
\end{tasks}

\ex{2} Primer: $\frac{18}{24}=75\,\%$. \ Segon: $\frac{22}{30}=73,\overline{3}\,\%$. \ \textbf{El primer} té millor percentatge.

\ex{2} Preu del viatge: $36:0,12=\res{300\text{ €}}$. \ Queda per pagar: $300-36=\res{264\text{ €}}$.

\ex{3}
\begin{tasks}(1)
  \task $\frac{162}{450}=0,36=\res{36\,\%}$.
  \task $40\,\%$ de $450=180$ hm$^3$. \ Ha guanyat $180-162=\res{18\text{ hm}^3}$.
\end{tasks}

\ex{3} $\frac38=0,375$ de la població; d'aquests, el $40\,\%$: \ $0,375\cdot0,4=0,15=\res{15\,\%}$ de la població total.

% =====================================================
\seccio{4}{Augments i disminucions percentuals}

\ex{1}\par\smallskip
{\renewcommand{\arraystretch}{1.6}\centering
\begin{tabular}{|C{3.2cm}|C{3.2cm}|C{3.2cm}|C{3.2cm}|}
\hline
\rowcolor{soft}\textbf{Quantitat inicial} & \textbf{Variació} & \textbf{Índex de variació} & \textbf{Quantitat final} \\ \hline
$200$ & augment del $15\,\%$ & $1,15$ & $230$ \\ \hline
$80$ & disminució del $25\,\%$ & $0,75$ & $60$ \\ \hline
$1\,200$ & augment del $3\,\%$ & $1,03$ & $1\,236$ \\ \hline
$500$ & disminució del $6\,\%$ & $0,94$ & $470$ \\ \hline
$64$ & augment del $25\,\%$ & $1,25$ & $80$ \\ \hline
\end{tabular}\par}

\ex{1}
\begin{tasks}(3)
  \task Augment del $12\,\%$
  \task Disminució del $15\,\%$
  \task Augment del $3\,\%$
  \task Disminució del $40\,\%$
  \task Augment del $150\,\%$
  \task Disminució del $0,5\,\%$
\end{tasks}

\ex{2} $640\cdot1,21=\res{774,40\text{ €}}$.

\ex{2} $750\cdot1,032=\res{774\text{ €}}$ al mes.

\ex{2} $I_V=\dfrac{11\,750}{12\,500}=0,94$ \ $\Rightarrow$ \ ha disminuït un $\res{6\,\%}$.

\ex{3} Primera botiga: $90\cdot1,1=99$ €. \ Segona botiga: $110\cdot0,9=99$ €. \ \textbf{Ara costa el mateix} en totes dues: $\res{99\text{ €}}$.

\ex{3} Dotze mesos: $45\cdot12=540$ €. \ Abonament anual: $540\cdot0,85=\res{459\text{ €}}$. \ S'estalvia $540-459=\res{81\text{ €}}$.

% =====================================================
\seccio{5}{Càlcul de la quantitat inicial}

\ex{1}
\begin{tasks}(4)
  \task $\dfrac{84}{1,2}=\res{70}$
  \task $\dfrac{56}{0,7}=\res{80}$
  \task $\dfrac{420}{1,05}=\res{400}$
  \task $\dfrac{220}{0,88}=\res{250}$
\end{tasks}

\ex{1} $\dfrac{181,50}{1,21}=\res{150\text{ €}}$.

\ex{2} Preu original: $\dfrac{520}{0,65}=\res{800\text{ €}}$. \ Estalvi: $800-520=\res{280\text{ €}}$.

\ex{2} $\dfrac{1\,802,50}{1,03}=\res{1\,750\text{ €}}$.

\ex{2} Un $40\,\%$ menys $\Rightarrow I_V=0,6$: \ $\dfrac{9\,600}{0,6}=\res{16\,000\text{ €}}$.

\ex{3}
\begin{tasks}(1)
  \task Sense IVA: $\dfrac{96,80}{1,21}=\res{80\text{ €}}$. \ IVA pagat: $96,80-80=\res{16,80\text{ €}}$ \ (que és el $21\,\%$ de $80$).
  \task L'IVA és el $21\,\%$ del preu \textbf{sense} IVA ($80$ €), no del preu final. Lluís aplica el percentatge a una quantitat que ja inclou l'IVA.
\end{tasks}

\ex{3} Tenia $\dfrac{3\,450}{0,92}=\res{3\,750\text{ L}}$, \ que és el $\dfrac{3\,750}{4\,000}=0,9375=\res{93,75\,\%}$ de la capacitat.

% =====================================================
\seccio{6}{Variacions percentuals encadenades}

\ex{1}
\begin{tasks}(2)
  \task $1,1\cdot1,2=1,32$ \ $\Rightarrow$ \ $\res{+32\,\%}$
  \task $0,9\cdot0,8=0,72$ \ $\Rightarrow$ \ $\res{-28\,\%}$
  \task $1,5\cdot0,5=0,75$ \ $\Rightarrow$ \ $\res{-25\,\%}$
  \task $1,25\cdot0,8=1$ \ $\Rightarrow$ \ $\res{0\,\%}$ (no varia)
\end{tasks}

\ex{1} $300\cdot1,05\cdot1,08=300\cdot1,134=\res{340,20\text{ €}}$.

\ex{2} $I_V=0,7\cdot0,9=0,63$. \ Paguem $120\cdot0,63=\res{75,60\text{ €}}$. \ Descompte total: $1-0,63=0,37$ \ $\Rightarrow$ \ $\res{37\,\%}$ (no el $40\,\%$).

\ex{2} $I_V=1,04\cdot0,94\cdot1,02=0,997152<1$. \ \textbf{Ha baixat} un $1-0,997152=0,002848$, és a dir, un $\res{0,28\,\%}$ aproximadament.

\ex{2} $5\,000\cdot1,12\cdot0,85=5\,000\cdot0,952=\res{4\,760\text{ €}}$. \ Hem \textbf{perdut} $240$ €, un $\res{4,8\,\%}$.

\ex{3} $I_V=\dfrac{1}{1,3}\approx0,7692$ \ $\Rightarrow$ \ ha de baixar aproximadament un $\res{23,08\,\%}$.

\ex{3} $1,2\cdot I_V=0,96$ \ $\Rightarrow$ \ $I_V=\dfrac{0,96}{1,2}=0,8$ \ $\Rightarrow$ \ ha baixat un $\res{20\,\%}$.

% =====================================================
\seccio{7}{Interés compost}

\ex{1} $C_f=3\,000\cdot1,04^{2}=\res{3\,244,80\text{ €}}$. \ Interessos: $\res{244,80\text{ €}}$.

\ex{1}\par\smallskip
{\renewcommand{\arraystretch}{1.6}\centering
\begin{tabular}{|C{2.2cm}|C{4cm}|C{4cm}|C{4cm}|}
\hline
\rowcolor{soft}\textbf{Any} & \textbf{Capital a l'inici} & \textbf{Interessos de l'any} & \textbf{Capital al final} \\ \hline
$1$ & $1\,000$ & $50$ & $1\,050$ \\ \hline
$2$ & $1\,050$ & $52,50$ & $1\,102,50$ \\ \hline
$3$ & $1\,102,50$ & $55,13$ & $1\,157,63$ \\ \hline
\end{tabular}\par}

\ex{2} $C_f=8\,000\cdot1,025^{5}\approx\res{9\,051,27\text{ €}}$ \ (interessos: $1\,051,27$ €).\\
Amb interés simple: $I=\dfrac{8\,000\cdot2,5\cdot5}{100}=\res{1\,000\text{ €}}$.

\ex{2} $C_0=\dfrac{10\,927,27}{1,03^{3}}=\dfrac{10\,927,27}{1,092727}\approx\res{10\,000\text{ €}}$.

\ex{2} Rèdit mensual: $3:12=0,25\,\%$; \ $24$ mesos: \ $C_f=12\,000\cdot1,0025^{24}\approx\res{12\,741,08\text{ €}}$.

\ex{3} Volem $1,04^{n}>\dfrac{6\,000}{5\,000}=1,2$. \ $1,04^{4}\approx1,1699$ ($5\,849,29$ €) i $1,04^{5}\approx1,2167$ ($6\,083,26$ €). \ Calen $\res{5\text{ anys}}$.

\ex{3} Banc A: $10\,000\cdot1,035^{4}\approx11\,475,23$ €. \ Banc B: $10\,000\cdot\left(1+\dfrac{3,4}{1\,200}\right)^{48}\approx11\,454,62$ €.\\
\textbf{És millor el banc A}, per $\res{20,61\text{ €}}$.

% =====================================================
\seccio{8}{Repartiments proporcionals}

\ex{1} $k=\dfrac{720}{3+4+5}=60$: \ $\res{180},\ \res{240},\ \res{300}$.

\ex{1} $k=\dfrac{1\,800}{10+15+20}=40$ €/h: \ $\res{400\text{ €}},\ \res{600\text{ €}},\ \res{800\text{ €}}$.

\ex{2} $\frac13,\frac14,\frac16=\frac{4}{12},\frac{3}{12},\frac{2}{12}$ \ $\Rightarrow$ \ directament a $4,3,2$: \ $k=\dfrac{1\,080}{9}=120$: \ $\res{480},\ \res{360},\ \res{240}$.

\ex{2} $\frac16,\frac19,\frac{1}{18}=\frac{3}{18},\frac{2}{18},\frac{1}{18}$ \ $\Rightarrow$ \ directament a $3,2,1$: \ $k=\dfrac{9\,000}{6}=1\,500$.\\
El de $6$ anys: $\res{4\,500\text{ €}}$; \ el de $9$ anys: $\res{3\,000\text{ €}}$; \ el de $18$ anys: $\res{1\,500\text{ €}}$.

\ex{2} $k=\dfrac{9\,600}{60\,000}=0,16$: \ $\res{2\,400\text{ €}},\ \res{3\,200\text{ €}},\ \res{4\,000\text{ €}}$.

\ex{3} Repartim directament a capital $\times$ temps: \ $48\,000$, $48\,000$ i $30\,000$ (suma $126\,000$).\\
$k=\dfrac{8\,400}{126\,000}=\dfrac{1}{15}$: \ A $\res{3\,200\text{ €}}$, \ B $\res{3\,200\text{ €}}$, \ C $\res{2\,000\text{ €}}$.

\ex{3} Repartim directament a $\dfrac{\text{hores}}{\text{faltes}}$: \ $\dfrac{30}{3}=10$ i $\dfrac{40}{2}=20$. \ $k=\dfrac{4\,500}{30}=150$: \ $\res{1\,500\text{ €}}$ i $\res{3\,000\text{ €}}$.

\ex{3} Les parts són $2k$, $3k$ i $5k$: \ $5k-2k=3k=450$ \ $\Rightarrow$ \ $k=150$.\\
Premi: $10\cdot150=\res{1\,500\text{ €}}$. \ Parts: $\res{300\text{ €}},\ \res{450\text{ €}},\ \res{750\text{ €}}$.

% =====================================================
\seccio{9}{Problemes de mescles}

\ex{1} $\dfrac{5\cdot6+3\cdot10}{5+3}=\dfrac{60}{8}=\res{7,50\text{ €/kg}}$.

\ex{1} $\dfrac{12\cdot2,5+8\cdot4}{20}=\dfrac{62}{20}=\res{3,10\text{ €/L}}$.

\ex{2} Balança: $20\cdot(7-6)=x\cdot(9-7)$ \ $\Rightarrow$ \ $20=2x$ \ $\Rightarrow$ \ $x=\res{10\text{ kg}}$.\\
(Comprovació: $\frac{20\cdot6+10\cdot9}{30}=\frac{210}{30}=7$ \faCheck)

\ex{2} Fruita: $3\cdot0,6+2\cdot0,2=2,2$ L de $5$ L: \ $\dfrac{2,2}{5}=0,44=\res{44\,\%}$.

\ex{2} Preu de cost: $\dfrac{30\cdot1,80+20\cdot2,40}{50}=\dfrac{102}{50}=2,04$ €/kg. \ Preu de venda: $2,04\cdot1,25=\res{2,55\text{ €/kg}}$.

\ex{3} Or pur: $40\cdot0,75+10=40$ g de $50$ g: \ llei $\dfrac{40}{50}=\res{0,800}$, \ és a dir, $0,8\cdot24=\res{19,2\text{ quirats}}$.

\ex{3} Balança (en quirats): l'or de $18$ k està $2$ k per davall de $20$ k i l'or pur ($24$ k), $4$ k per damunt:\\
$60\cdot2=x\cdot4$ \ $\Rightarrow$ \ $x=\res{30\text{ g}}$. \ (Comprovació: $\frac{45+30}{90}=\frac56=\frac{20}{24}$ \faCheck)

\ex{3} Alcohol pur: $15\cdot0,7=10,5$ L, que ha de ser el $40\,\%$ del total: \ $10,5:0,4=26,25$ L.\\
Cal afegir $26,25-15=\res{11,25\text{ L}}$ d'aigua.

% =====================================================
\seccio{10}{Problemes de moviments}

\ex{1} $54:3,6=\res{15\text{ m/s}}$. \quad $25\cdot3,6=\res{90\text{ km/h}}$.

\ex{1}
\begin{tasks}(1)
  \task $t=\dfrac{330}{120}=2,75$ h $=\res{2\text{ h }45\text{ min}}$.
  \task $1$ h $20$ min $=\frac43$ h: \ $d=120\cdot\frac43=\res{160\text{ km}}$.
\end{tasks}

\ex{2} S'acosten a $70+90=160$ km/h: \ $t=\dfrac{360}{160}=2,25$ h $=2$ h $15$ min \ $\Rightarrow$ \ a les $\res{11{:}15}$.\\
Distància a la ciutat del primer cotxe: $70\cdot2,25=\res{157,5\text{ km}}$.

\ex{2} Avantatge: $5\cdot2=10$ km. \ $t=\dfrac{10}{15-5}=\res{1\text{ h}}$, \ a $15\cdot1=\res{15\text{ km}}$ del punt de partida.

\ex{2} S'allunyen a $18+22=40$ km/h: \ $t=\dfrac{100}{40}=\res{2,5\text{ h}}=2$ h $30$ min.

\ex{3} A les $9{:}00$ el camió ha fet $80\cdot1,5=120$ km; \ els separen $500-120=380$ km.\\
$t=\dfrac{380}{80+110}=2$ h \ $\Rightarrow$ \ es creuen a les $\res{11{:}00}$, a $80\cdot3,5=\res{280\text{ km}}$ d'A \ (el cotxe ha fet $220$ km).

\ex{3} Avantatge: $60\cdot1,5=90$ km. \ $t=\dfrac{90}{150-60}=\res{1\text{ h}}$, \ a $150\cdot1=\res{150\text{ km}}$ de l'estació.

\ex{3} Ha de traure-li $400$ m i en guanya $1$ cada segon: \ $t=\res{400\text{ s}}=6$ min $40$ s.\\
La ràpida fa $6\cdot400=2\,400$ m ($\res{6\text{ voltes}}$) i l'altra $2\,000$ m ($\res{5\text{ voltes}}$).

% =====================================================
\seccio{11}{Proporcionalitat composta}

\ex{1}
\begin{tasks}(2)
  \task Inversament proporcional.
  \task Directament proporcional.
  \task No proporcional.
  \task Inversament proporcional.
  \task No proporcional (àrea $=\ell^{2}$).
  \task Directament proporcional.
\end{tasks}

\ex{1}
\begin{tasks}(1)
  \task Directa: $\dfrac{26}{4}\cdot7=\res{45,50\text{ €}}$.
  \task Inversa: $\dfrac{8\cdot15}{12}=\res{10\text{ dies}}$.
\end{tasks}

\ex{2} Aixetes (D) i hores (D): \ $x=3\,600\cdot\dfrac96\cdot\dfrac54=\res{6\,750\text{ L}}$.

\ex{2} Vaques (D) i dies (D): \ $x=1\,800\cdot\dfrac{15}{10}\cdot\dfrac{20}{12}=\res{4\,500\text{ kg}}$.

\ex{2} Obrers (I) i hores/dia (I): \ $x=18\cdot\dfrac56\cdot\dfrac{8}{10}=\res{12\text{ dies}}$.

\ex{3} Impressores (I), hores/dia (I) i fulls (D):\\
$x=5\cdot\dfrac{12}{10}\cdot\dfrac68\cdot\dfrac{72\,000}{54\,000}=\dfrac{5\cdot12\cdot6\cdot72\,000}{10\cdot8\cdot54\,000}=\res{6\text{ dies}}$.

\ex{3} Pinso (D) i dies (I): \ $x=30\cdot\dfrac{90}{120}\cdot\dfrac{20}{25}=\res{18\text{ gallines}}$.

\ex{3} Obrers (I), hores/dia (I) i metres (D):\\
$x=15\cdot\dfrac{8}{10}\cdot\dfrac96\cdot\dfrac{200}{120}=\dfrac{15\cdot8\cdot9\cdot200}{10\cdot6\cdot120}=\res{30\text{ dies}}$.

\end{document}
